\section{The elephant walk and its crossing}\label{sec:erw}

We reduce the elephant walk to a counting chain, prove that the crossing is
unique, compute the center to first and second order in $\varepsilon$, and
locate the crossing. The last subsection gives one more term of the center
and of the crossing.

\begin{definition}[elephant walk]\label{def:erw}
Fix $0<\varepsilon<1$. Let $X_1$ be uniform on $\{\pm1\}$. For $t\ge1$ let
$U_t$ be uniform on $\{1,\dots,t\}$, and put $X_{t+1}=X_{U_t}$ with
probability $1-\varepsilon$ and $X_{t+1}=-X_{U_t}$ with probability
$\varepsilon$. All these choices are independent. We call the event
$X_{t+1}=-X_{U_t}$ a flip. As in Section~\ref{sec:notation},
$A_N=\#\{i\le N:X_i=+1\}$ and $P_N(k)=\Prob(A_N=k)$. The crossing
$\varepsilon_N$ is the value of $\varepsilon$ with $C_N=E_N$, and
$x_N=N\varepsilon_N$.
\end{definition}

A flip is what the introduction called a reversal. In the notation
of~\cite{schutz2004} the memory parameter is $1-\varepsilon$,
and~\cite{glrs2025} write $a=1-2\varepsilon$. Given the first step, all that
matters is how many steps disagree with it. So let $M_t$ be the number of
$i\le t$ with $X_i\ne X_1$, and put
\[
 r_t(b)=\varepsilon+(1-2\varepsilon)\frac bt,\qquad f_N(m)=\Prob_+(M_N=m).
\]

\begin{lemma}\label{lem:erw-reduction}
The process $(M_t)_{t\ge1}$ is a Markov chain with $M_1=0$ and
$\Prob(M_{t+1}=b+1\mid M_t=b)=r_t(b)=1-\Prob(M_{t+1}=b\mid M_t=b)$, and it
is independent of $X_1$. For $0\le k\le N$,
$P_N(k)=\frac12[f_N(N-k)+f_N(k)]$. In particular $C_N=f_N(h)$,
$E_N=\frac12(1-\varepsilon)^{N-1}$, and $P_N(k)=P_N(N-k)$.
\end{lemma}

\begin{proof}
Given $X_1,\dots,X_t$, the copied step $X_{U_t}$ differs from $X_1$ with
probability $M_t/t$. So $X_{t+1}\ne X_1$ with conditional probability
$(1-\varepsilon)M_t/t+\varepsilon(1-M_t/t)=r_t(M_t)$. This depends only on
$M_t$ and not on $X_1$, which proves the first claim. On $\{X_1=+1\}$ we
have $A_N=N-M_N$, and on $\{X_1=-1\}$ we have $A_N=M_N$. Averaging over
$X_1$ gives the formula for $P_N(k)$, which does not change under
$k\mapsto N-k$. Finally $M_N\le N-1$, so $f_N(N)=0$, and
$f_N(0)=\prod_{t=1}^{N-1}(1-r_t(0))=(1-\varepsilon)^{N-1}$.
\end{proof}

So the end is Kagey's end $p^{N-1}/2$ with $p=1-\varepsilon$. More flips
always help the middle against the end.

\begin{proposition}\label{prop:mlr}
Let $N\ge2$ and $1\le k\le N-1$. Then $f_N(k)/f_N(0)$ and $P_N(k)/E_N$ are
strictly increasing in $\varepsilon\in(0,1)$. In particular $C_N/E_N$ is
strictly increasing, the crossing $\varepsilon_N$ exists and is unique, and
$\varepsilon_N<\frac12$. Also $C_N\ge2E_N$ for $\varepsilon\ge\frac12$.
\end{proposition}

\begin{proof}
The probability of a path of $M$ is a product of $N-1$ factors, namely
$r_s(b_s)$ for an up step at time $s$ and $1-r_s(b_s)$ otherwise, where $b_s$
is the state at time $s$. We divide it by the probability
$f_N(0)=(1-\varepsilon)^{N-1}$ of the path that never flips. With
$y=b_s/s\in[0,1)$ the factors become
$\frac{\varepsilon+(1-2\varepsilon)y}{1-\varepsilon}$ and
$1-y\frac{1-2\varepsilon}{1-\varepsilon}$, whose derivatives in
$\varepsilon$ are $\frac{1-y}{(1-\varepsilon)^2}>0$ and
$\frac y{(1-\varepsilon)^2}\ge0$. Both factors are positive, since
$r_s(b_s)=\varepsilon(1-y)+(1-\varepsilon)y$ lies strictly between 0 and 1.
So the ratio for each path is nondecreasing in $\varepsilon$, and strictly
increasing if the path has an up step. Summing over the paths with $k\ge1$
up steps proves the claim for $f_N(k)/f_N(0)$, and the claim for $P_N(k)/E_N$
follows by Lemma~\ref{lem:erw-reduction}. As $\varepsilon\to0$, $C_N\to0$
and $E_N\to\frac12$. For
$\varepsilon\ge\frac12$ every $r_s(b)$ and every $1-r_s(b)$ lies in
$[1-\varepsilon,\varepsilon]$, so each path of $M$ that ends at $h$ has
probability at least $(1-\varepsilon)^{N-1}=2E_N$. Hence $C_N\ge2E_N$. The
ratio $C_N/E_N$ is continuous and strictly increasing, so it equals 1 at
exactly one $\varepsilon$, and this $\varepsilon$ lies in $(0,\frac12)$.
\end{proof}

\subsection{First order, one early flip}

When $\varepsilon=0$ the chain $M$ never leaves 0. After a first flip, and to
first order in $\varepsilon$, it is P\'olya's urn.

\begin{lemma}[P\'olya's urn]\label{lem:urn}
Let $2\le v\le N$. At time $v$ an urn holds one ball of one color and $v-1$
of the other. At each time $k=v,\dots,N-1$ we draw a ball uniformly and
return it with one more ball of its color. Let $\beta_v(d)$ be the
probability that at time $N$ there are $d$ balls of the first color. Then
$\beta_v(d)=\binom{N-d-1}{v-2}/\binom{N-1}{v-1}$ for $1\le d\le N-v+1$, and
$\beta_v(d)=0$ for other $d$. Also $\beta_v$ is nonincreasing in $d$ and
$\sum_dd\,\beta_v(d)=N/v$. For $1\le d\le N-1$,
\[
 \sum_{v=2}^N\beta_v(d)=\frac{N}{d(d+1)},\qquad
 \sum_{v=2}^N(v-2)\,\beta_v(d)=\frac{2N(N-d-1)}{d(d+1)(d+2)} .
\]
\end{lemma}

\begin{proof}
If the urn holds $s$ balls of the first color at time $k$, the next draw
is of that color with probability $s/k$. So a given order of $i$ draws of
the first color among the $N-v$ draws has probability
$\frac{i!\,(v-1)v\cdots(N-i-2)}{v(v+1)\cdots(N-1)}=\frac{(v-1)\,i!\,(N-i-2)!}{(N-1)!}$.
Multiplying by $\binom{N-v}{i}$ and putting $d=i+1$ gives
$\beta_v(d)=a_d(v-1)\binom{N-v}{d-1}$ with $a_d=\frac{(d-1)!(N-d-1)!}{(N-1)!}$,
which equals $\binom{N-d-1}{v-2}/\binom{N-1}{v-1}$. Hence
$\beta_v(d)/\beta_v(d-1)=1-\frac{v-2}{N-d}$ for $2\le d\le N-v+1$, and
$\beta_v$ is nonincreasing. The fraction of balls of the first color is a
martingale that starts at $1/v$, which gives the mean. For the two sums put
$i=N-v$, so that $v-1=\binom{N-1-i}1$ and $(v-1)(v-2)=2\binom{N-1-i}2$. By
the Chu--Vandermonde identity,
$\sum_i\binom{i}{d-1}\binom{N-1-i}{1}=\binom{N}{d+1}$ and
$\sum_i\binom{i}{d-1}\binom{N-1-i}{2}=\binom{N}{d+2}$. Multiplying by $a_d$
and $2a_d$ gives the two sums.
\end{proof}

We need some more notation. Let $V_s(b)=\Prob_+(M_N=h\mid M_s=b)$ be the
chance of ending at the center from state $b$ at time $s$. For $1\le t\le h$
let
\[
 \pi_t=\beta_{t+1}(h)=\frac t{N-t}\prod_{i=1}^{t-1}\frac{h-i}{N-i},\qquad
 c_1=\sum_{t=1}^h\pi_t,\qquad w_t=\frac{\pi_t}{c_1},
\]
so that $(w_t)_{1\le t\le h}$ is a probability vector.

\begin{proposition}\label{prop:first-mutation}
Let $N\ge2$ be even.
\begin{enumerate}[(a)]
\item $C_N=\sum_{t=1}^h\varepsilon(1-\varepsilon)^{t-1}V_{t+1}(1)$.
\item $C_N$ is a polynomial in $\varepsilon$ with $C_N(0)=0$. At
$\varepsilon=0$ we have $V_{t+1}(1)=\pi_t$, and
$C_N'(0)=c_1=\frac4{N+2}$ exactly.
\item $\sum_tw_tt=3-\frac6{h+2}$, and $w_t\le2t2^{-t}$ for $1\le t\le h$.
\end{enumerate}
\end{proposition}

\begin{proof}
(a) Under $\Prob_+$ the first flip happens at step $t+1$ exactly when
$M_s=0$ for $s\le t$ and $M_{t+1}=1$. Since $r_s(0)=\varepsilon$, this has
probability $(1-\varepsilon)^{t-1}\varepsilon$, and by the Markov property
the chain then ends at $h$ with probability $V_{t+1}(1)$. At time $t+1$
there are $t$ steps equal to $X_1$, and this number never decreases. So
$V_{t+1}(1)=0$ for $t>h$. Paths with no flip end at $0\ne h$.

(b) Every transition probability is affine in $\varepsilon$, so $C_N$ and
$V_{t+1}(1)$ are polynomials in $\varepsilon$. By (a), $C_N(0)=0$ and
$C_N'(0)=\sum_tV_{t+1}(1)|_{\varepsilon=0}$. At $\varepsilon=0$ we have
$r_s(b)=b/s$. So from time $t+1$ on, $M$ counts the balls of the first
color in the urn of Lemma~\ref{lem:urn} with $v=t+1$, and
$V_{t+1}(1)=\beta_{t+1}(h)=\pi_t$. Since $\beta_v(h)=0$ for $v>h+1$,
Lemma~\ref{lem:urn} with $d=h$ gives
$c_1=\sum_{v=2}^N\beta_v(h)=\frac N{h(h+1)}=\frac4{N+2}$.

(c) In the same way the second sum of Lemma~\ref{lem:urn} gives
$\sum_t(t-1)\pi_t=\frac{2N(h-1)}{h(h+1)(h+2)}$, so
$\sum_tw_t(t-1)=\frac{2(h-1)}{h+2}$. Next $\frac{h-i}{N-i}\le\frac12$, so
$w_t=\frac{N+2}4\pi_t\le\frac{N+2}{2(N-t)}\,t2^{-t}\le2t2^{-t}$ for $t\le h$.
\end{proof}

In words, to first order the walk flips once, at some step $t+1$, and then
runs as P\'olya's urn from $t$ balls of one color and one of the other. The
minority fraction is then close to a Beta$(1,t)$ variable, whose density at
$\frac12$ is $t2^{1-t}$, and $\sum_tt2^{1-t}=4$ explains the constant
$4/(N+2)$. Since $w_t\le2t2^{-t}$, only an early flip counts.

\subsection{Second order}

\begin{theorem}\label{thm:erw-bounds}
Let $N\ge2$ be even.
\begin{enumerate}[(a)]
\item For $0<\varepsilon<\frac12$,
$C_N\ge\frac{4\varepsilon}{N+2}(1-2\varepsilon)$.
\item $C_N\le\frac{4\varepsilon}{N+2}\,N^{2\varepsilon}e^{5200\varepsilon}$
whenever $N\ge64$ and $0<\varepsilon(1+\log N)\le\frac1{100}$.
\end{enumerate}
In particular $\log C_N=\log\frac{4\varepsilon}{N+2}+O(\varepsilon\log N)$
in the range of (b), with an absolute implied constant.
\end{theorem}

In words, the first-order term is correct up to a factor
$1+O(\varepsilon\log N)$, which is $1+O(L^2/N)$ at the crossing.

\paragraph{Idea of the proof.}
The proof is in Appendix~\ref{app:erw}. After the first flip, $M$ at time $s$
is P\'olya's urn with $\beta s$ extra balls of each color, where
$\beta=\varepsilon/(1-2\varepsilon)$ (Lemma~\ref{lem:A-polya}). They pull the
minority fraction toward $\frac12$. For (a) we drop them. For (b) we compare
$M$ with an urn with $\kappa(s+10)$ extra balls, where
$\kappa=2\varepsilon/(1-4\varepsilon)$. The step at time $s$ then loses a
factor of about $e^{\kappa/s}$, and the product of these factors is about
$e^{\kappa H_N}\approx N^{2\varepsilon}$.

\subsection{The crossing}

By Theorem~\ref{thm:erw-bounds},
$C_N\approx4x/N^2$ when $\varepsilon=x/N$ and $x$ is of order $\log N$. Also
$E_N=\frac12(1-x/N)^{N-1}\approx\frac12e^{-x}$. Setting them equal gives
$x+\log x\approx2\log N-\log8$. So let
\[
 \Lambda_N=2L-\log8,\qquad G(x)=x+\log x,
\]
and let $g_N$ be the root of $G(g_N)=\Lambda_N$. It exists and is unique,
since $G$ increases from $-\infty$ to $\infty$ on $(0,\infty)$.

\begin{table}[t]
\centering\small
\begin{tabular}{@{}rlllll@{}}
\toprule
$N$ & $x_N$ & $g_N$ & $x_N-x_{\mathrm{2nd}}$ & $x_N/\log N$ & $D(N)\log N$\\
\midrule
50 & 3.940487 & 4.288636 & $+0.3452$ & 1.007 & $-0.083$\\
400 & 7.629217 & 7.843768 & $+0.0621$ & 1.273 & $-0.199$\\
1600 & 10.240565 & 10.340051 & $+0.0179$ & 1.388 & $+0.215$\\
3200 & 11.546834 & 11.610464 & $+0.0097$ & 1.431 & $+0.386$\\
12800 & 14.158845 & 14.182921 & $+0.0028$ & 1.497 & $+0.607$\\
51200 & 16.778468 & 16.786946 & $+0.0008$ & 1.547 & $+0.707$\\
\bottomrule
\end{tabular}
\caption{Exact elephant crossings, the root $g_N$ of $G(g)=\Lambda_N$, and
the difference $D(N)$ of Theorem~\ref{thm:erw-markov} times $\log N$. Here
$x_{\mathrm{2nd}}$ is the formula of Proposition~\ref{prop:refined} without
its error term.}
\label{tab:erw-crossing}
\end{table}

\begin{theorem}\label{thm:erw-crossing}
As $N\to\infty$ over even $N$, the following hold.
\begin{enumerate}[(a)]
\item $x_N+\log x_N=\Lambda_N+O(L^2/N)$, so $x_N=g_N+O(L^2/N)$.
\item $x_N=2\log N-\log\log N-\log16+O(\log\log N/\log N)$.
\end{enumerate}
\end{theorem}

In words, the root $g_N$ locates the crossing to within $O(L^2/N)$, and the
expected number of flips at the crossing is about $2\log N$.

\paragraph{Idea of the proof.}
Let $F(x)$ be the logarithm of $C_N/E_N$ at $\varepsilon=x/N$. It increases
in $x$, and its only zero is $x_N$. Theorem~\ref{thm:erw-bounds} and the
formula for $E_N$ show that $F(x)$ is within $O(L^2/N)$ of $G(x)-\Lambda_N$
when $x$ is of order $L$. Since $G'\ge1$, the zero of $F$ is within
$O(L^2/N)$ of the zero $g_N$ of $G-\Lambda_N$.

\begin{proof}
(a) \emph{Step 1: both sides near the crossing.} Put
$F(x)=\log C_N(x/N)-\log E_N(x/N)$ for $0<x<N$, where $C_N(\varepsilon)$ and
$E_N(\varepsilon)$ are the center and the end at flip probability
$\varepsilon$. By Proposition~\ref{prop:mlr}, $F$ is strictly increasing and
$x_N$ is its only zero. Let $1\le x\le2\Lambda_N$. Then $x/N\le4L/N$, so for
large $N$ Theorem~\ref{thm:erw-bounds} applies with $\varepsilon=x/N$. As
$\log\frac{4\varepsilon}{N+2}=\log x+\log4-2L-\log(1+2/N)$ and $x\ge1$, it
gives $\log C_N=\log x+\log4-2L+O(xL/N)$. Also $\log(1-y)=-y+O(y^2)$ for
$0\le y\le\frac12$, so $\log E_N=-\log2-(N-1)x/N+O(x^2/N)=-\log2-x+O(x^2/N)$.
Since $\log4+\log2=\log8$ and $x\le2\Lambda_N\le4L$, we get
$|F(x)-G(x)+\Lambda_N|\le\rho_N$ on $[1,2\Lambda_N]$, with $\rho_N=O(L^2/N)$.

\emph{Step 2: trap the zero.} Now $G'\ge1$ and $G(g_N)=\Lambda_N$. So
$F(x)<0$ for $1\le x<g_N-\rho_N$ and $F(x)>0$ for $g_N+\rho_N<x\le2\Lambda_N$.
As $\Lambda_N-\log\Lambda_N\le g_N\le\Lambda_N$ (because
$G(\Lambda_N-\log\Lambda_N)\le\Lambda_N\le G(\Lambda_N)$), both ranges are
nonempty for large $N$. Since $F$ is increasing, its zero $x_N$ lies between
them, so $|x_N-g_N|\le\rho_N$. Then
$F(x_N)=0$ gives $|G(x_N)-\Lambda_N|\le\rho_N$.

(b) Since $g_N=\Lambda_N-\log g_N$ and
$\Lambda_N-\log\Lambda_N\le g_N\le\Lambda_N$, we have
$\log g_N=\log\Lambda_N+O(\log\Lambda_N/\Lambda_N)$, and
$\log\Lambda_N=\log2+\log\log N+O(1/\log N)$. So
$g_N=2\log N-\log8-\log2-\log\log N+O(\log\log N/\log N)$, which is (b) for
$g_N$, and (a) gives it for $x_N$.
\end{proof}

Table~\ref{tab:erw-crossing} gives exact crossings on the grid
$N=50\cdot2^k$, computed from the recursion for $f_N$ in
Lemma~\ref{lem:erw-reduction}. The limit $2$ of $x_N/\log N$ is approached
slowly, because of the term $-\log\log N$.

Now we compare with Kagey's board. Doubling \eqref{eq:markov-crossing}
without its last two terms gives $2z_N+\log(2z_N)=\Lambda_N+\log(2\pi)$, so
$2z_N-\log(2\pi)$ nearly solves $G(x)=\Lambda_N$. What is left is of order
$1/L$.

\begin{theorem}\label{thm:erw-markov}
Let $z_N$ be the crossing value of $T$ for Kagey's board with $N$ steps, as
in Section~\ref{sec:series}. As $N\to\infty$ over even $N$,
\[
 D(N):=x_N-\bigl(2z_N-\log(2\pi)\bigr)=\frac{c_*}{\log N}
 +O\Bigl(\frac{\log\log N}{(\log N)^2}\Bigr),\qquad
 c_*=\tfrac12\log(2\pi)-\tfrac14=0.668939\ldots
\]
\end{theorem}

In words, the elephant needs twice as many flips as Kagey's board needs
switches, less $\log(2\pi)=1.837\ldots$, and the difference from this tends
to 0 like $c_*/\log N$.

\paragraph{Idea of the proof.}
Both crossings solve an equation of the form $x+\log x=2L+\text{constant}$
up to small errors. We put the Markov crossing into the equation
$G(x)=\Lambda_N$ of the elephant, find that $2z_N-\log(2\pi)$ misses it by
$c_*/L$, and divide by $G'\approx1$.

\begin{proof}
Write $\ell=\log(2\pi)$. By Section~\ref{sec:series}, the error term in
\eqref{eq:markov-crossing} is $O(L^{-2})$ and $L/2\le z_N\le L$. So
\eqref{eq:markov-crossing} gives $z_N+\frac12\log z_N=L+O(1)$ with
$\log z_N=\log L+O(1)$, hence $z_N=L-\frac12\log L+O(1)$. Since $z_N\ge L/2$,
this gives $\frac1{8z_N}-\frac1{8L}=\frac{L-z_N}{8Lz_N}=O(\log L/L^2)$. Put
$Y=2z_N$. Twice \eqref{eq:markov-crossing} plus $\log2$, together with
$2c_0+\log2=\log\frac\pi4$, gives
$Y+\log Y=2L+\log\frac\pi4+\frac1{4L}+O(\log L/L^2)$. Also
$Y=2L-\log L+O(1)$, so $1/Y=1/(2L)+O(\log L/L^2)$ and
$\log(1-\ell/Y)=-\ell/(2L)+O(\log L/L^2)$. Using $\log\frac\pi4-\ell=-\log8$,
we get
\[
 G(Y-\ell)=Y+\log Y-\ell+\log\Bigl(1-\frac\ell Y\Bigr)
 =\Lambda_N+\frac1{4L}-\frac\ell{2L}+O\Bigl(\frac{\log L}{L^2}\Bigr)
 =\Lambda_N-\frac{c_*}L+O\Bigl(\frac{\log L}{L^2}\Bigr).
\]
By Theorem~\ref{thm:erw-crossing}(a), $G(x_N)=\Lambda_N+O(L^2/N)$. Both
$x_N$ and $Y-\ell=2z_N-\log(2\pi)$ are $2L+O(\log L)$, so $G'=1+1/x=1+O(1/L)$
between them, and the mean value theorem gives
$x_N-(Y-\ell)=c_*/L+O(\log L/L^2)$.
\end{proof}

We computed $z_N$ from the exact law of Kagey's board~\cite{paperA}. The
difference $D(N)$ changes sign between $N=400$ and $N=1600$, and
$D(N)\log N$ is $0.707$ at $N=51200$ against the limit $c_*=0.669$
(Table~\ref{tab:erw-crossing}). So the asymptotic regime is reached slowly.

\subsection{The crossing to second order}\label{sec:erw-second}

Theorem~\ref{thm:erw-bounds} gives $\log C_N$ up to an error
$O(\varepsilon\log N)$. A lower bound that matches part (b) of that theorem
brings the error down to $O(\varepsilon)$.

\begin{proposition}[the center to second order]\label{prop:second-order}
Let $N\ge64$ be even and $0<\varepsilon(1+\log N)\le\frac1{100}$. Then
\[
 -11\varepsilon-\varepsilon^2(84\log N+1.05\cdot10^5)
 \le\log C_N-\log\frac{4\varepsilon}{N+2}-2\varepsilon\log N\le5200\,\varepsilon .
\]
In particular
$\bigl\lvert\log C_N-\log\frac{4\varepsilon}{N+2}-2\varepsilon\log N\bigr\rvert
\le5200\,\varepsilon$ there.
\end{proposition}

In words, $C_N$ is $\frac{4\varepsilon}{N+2}N^{2\varepsilon}$ times a factor
$1+O(\varepsilon)$.

\paragraph{Idea of the proof.}
The proof is in Appendix~\ref{app:erw}. The upper bound is
Theorem~\ref{thm:erw-bounds}(b). For the lower bound we compare $M$ after
the first flip with P\'olya's urn conditioned to end at $h$, which we call
the bridge. Then $V_{t+1}(1)/\pi_t$ is the mean over the bridge of a product
of factors $1+O(\varepsilon)$, one for each step
(Proposition~\ref{prop:A-bridge}). By Jensen's inequality the mean of the
product is at least the exponential of the mean of the sum of the
logarithms. The step at time $u$ adds about $2\varepsilon/u$ to this mean,
and the sum over $u\le N$ is about $2\varepsilon\log N$.

With this bound the proof of Theorem~\ref{thm:erw-crossing} gives one more
term of the crossing.

\begin{proposition}[the crossing to second order]\label{prop:refined}
As $N\to\infty$ over even $N$,
\[
 x_N=g_N-\frac{2g_N\log N+g_N^2/2}{N(1+1/g_N)}+O\Bigl(\frac{\log N}N\Bigr).
\]
So $x_N<g_N$ for large $N$.
\end{proposition}

In words, the root $g_N$ overshoots the crossing by about $6\log^2N/N$.

\begin{proof}
Let $F$ be as in the proof of Theorem~\ref{thm:erw-crossing}, and let
$1\le x\le2\Lambda_N$ and $\varepsilon=x/N$. For large $N$
Proposition~\ref{prop:second-order} applies and gives
$\log C_N=\log x+\log4-2L+\frac{2xL}N+O(\frac{x}N)$. For
$0\le y\le\frac12$ we have $\log(1-y)=-y-\frac{y^2}2+O(y^3)$, so
$\log E_N=-\log2+(N-1)\log(1-\frac xN)=-\log2-x-\frac{x^2}{2N}+O(\frac xN)$,
because $x^2/N^2$ and $x^3/N^2$ are $O(x/N)$ for $x\le4L$. Hence
\[
 F(x)=G(x)-\Lambda_N+\frac{2xL+x^2/2}{N}+O\Bigl(\frac xN\Bigr)
 \qquad(1\le x\le2\Lambda_N).
\]
By Theorem~\ref{thm:erw-crossing}, $x_N$ lies in this range for large $N$.
So $F(x_N)=0$ gives $G(x_N)=\Lambda_N-\eta_N$ with
$\eta_N=(2x_NL+x_N^2/2)/N+O(L/N)$. By the mean value theorem,
$x_N-g_N=-\eta_N/G'(\xi)$ for some $\xi$ between $x_N$ and $g_N$. Here
$G'(\xi)=1+1/\xi=(1+1/g_N)(1+O(1/N))$, since $\lvert\xi-g_N\rvert=O(L^2/N)$
and $g_N\ge L$ for large $N$. As $\eta_N=O(L^2/N)$, we get
$x_N-g_N=-\eta_N/(1+1/g_N)+O(L^2/N^2)$. Replacing $x_N$ by $g_N$ in the main
term of $\eta_N$ changes it by $O(L^3/N^2)$. This gives the formula. Its
main term is of order $L^2/N$ and is larger than the error for large $N$, so
$x_N<g_N$.
\end{proof}

Proposition~\ref{prop:second-order} does not identify the term of order
$\varepsilon$. For fixed $N$ and small $\varepsilon$ we can find it. By
Proposition~\ref{prop:first-mutation}(b), $C_N$ is a polynomial in
$\varepsilon$, and we write
$C_N=c_1\varepsilon+c_2\varepsilon^2+O(\varepsilon^3)$ as $\varepsilon\to0$
and $k_2(N)=c_2/c_1$.

\begin{proposition}[the coefficient of $\varepsilon^2$]\label{prop:K2}
As $N\to\infty$ over even $N$,
\[
 k_2(N)-2\log N\to2\gamma-2-\log2=-1.5387\ldots
\]
\end{proposition}

In words, for fixed large $N$ and small $\varepsilon$ the factor
$1+O(\varepsilon)$ in Proposition~\ref{prop:second-order} is about
$1-1.54\,\varepsilon$.

\paragraph{Idea of the proof.}
The proof is in Appendix~\ref{app:erw}. Differentiating the bridge formula
at $\varepsilon=0$ gives $k_2(N)=\sum_tw_t[A_t-(t-1)]$, where $A_t$ is the
mean over the bridge of an explicit sum (Lemma~\ref{lem:A-K2}). For fixed
$t$ and large $N$ the bridge is close to a walk with fair steps, and the
sum can be computed in closed form. It is
$2\log N+2(\gamma-1-\log2)+(t-2)H_{t-1}+o(1)$. The weights $w_t$ tend to
$t2^{-t-1}$, and the series over $t$ sums to $2\gamma-2-\log2$.

\begin{remark}[numerical values and open points]\label{rem:refined}
From the Taylor coefficients of $C_N$, computed in double precision
(\texttt{verify\_elephant\_crossing.py}), we get
$k_2(N)-2\log N=-1.585878$, $-1.546572$ and $-1.541005$ at $N=1600$, $12800$
and $51200$. The formula of Proposition~\ref{prop:refined} without its error
term is within $0.07$ of $x_N$ already at $N=400$
(Table~\ref{tab:erw-crossing}). We do not know an upper bound
$\log C_N\le\log(c_1\varepsilon)+\varepsilon k_2(N)+O(\varepsilon^2(\log N)^k)$,
or an explicit $N$ beyond which $g_N>x_N$.
\end{remark}
